\documentclass[twocolumn,10pt]{ctexart}
\usepackage[UTF8,fontset=windows]{ctex}
\begin{document}
\twocolumn[
\begin{center}
{\heiti\zihao{3} FAECO：面向预布局逻辑重综合的失效驱动候选生成引擎}\par
\vspace{6pt}
{\fangsong\zihao{4} 佟亚龙}\par
\vspace{4pt}
{\songti\fontsize{8}{12}\selectfont（国防科技大学计算机学院，长沙 410073）}\par\vspace{2pt}
{\songti\fontsize{8}{12}\selectfont E-mail：tyl2003@nudt.edu.cn（通信作者）}\par
\vspace{8pt}
\end{center}
\begin{flushleft}
{\heiti\fontsize{8.5}{13}\selectfont 摘\quad 要：}{\songti\fontsize{8.5}{13}\selectfont 工程变更命令修复中，缓冲器插入与门尺寸调整仅能微调门级延迟，对逻辑级数过深的关键路径无能为力；逻辑重综合虽能改变电路结构，但其候选构造阶段将等价性验证与物理时序优化串行处理，导致「验证失败」与「收益不足」两类问题相互制约。为此，提出面向预布局逻辑重综合的失效驱动候选生成引擎 FAECO：将等价性与时序收益双目标编码为割图权重偏好，按六类失败模式自适应更新权重并重新搜索割解，并将估计线载嵌入内环接受准则，以理想线网改善阈值与简化 SPEF 复测两阶段门控筛选物理稳健候选。在 ISCAS89（8/8）、ITC-99（18/19）及 PicoRV32（2/3）基准上取得严格时序改善；收敛配置下平均 WNS 改善 1.07 ns（各电路提升幅度的均值），为纯门尺寸调整的 6.7 倍；效率优先模式下仅需 27 次 STA 调用即完成 8/8 修复，开销削减 99.1%。经真实布局布线验证，8 个电路中 5 个保持改善、1 个持平、2 个回退不超过 0.02 ns。FAECO 在开源工具链上实现了失效驱动逻辑重综合的闭环修复。}
\par
\vspace{4pt}
{\heiti\fontsize{8.5}{13}\selectfont 关键词：}{\songti\fontsize{8.5}{13}\selectfont 时序 ECO；重综合；双目标割；失效驱动反馈；简化 SPEF 门控；预布局逻辑重综合}
\vspace{2pt}
{\songti\fontsize{8.5}{13}\selectfont 中图法分类号：TP391.41\qquad DOI：10.3724/SP.J.1089.2026-00001}
\end{flushleft}
\vspace{6pt}
% 英文题录
\begin{center}
{\bfseries\fontsize{13.5}{18}\selectfont FAECO: Failure-Aware ECO Candidate Generation Engine for Pre-Placement Logic Re-Synthesis}\par
\vspace{5pt}
{\normalsize Tong Yalong\par}
\vspace{3pt}
{\small \itshape (College of Computer Science, National University of Defense Technology, Changsha 410073, China)}\par
\vspace{6pt}
\end{center}
\begin{flushleft}
{\bfseries Abstract:} {\small In engineering change order (ECO) repair, buffer insertion and gate sizing can only tune gate-level delays and cannot handle critical paths with excessive logic depth; logic re-synthesis can change the circuit structure, but its candidate construction stage processes equivalence checking and physical timing optimization serially, causing the mutual restriction of verification failures and insufficient timing gains. To address this, FAECO (Failure-Aware ECO), a failure-aware candidate generation engine for pre-placement logic re-synthesis, is proposed. It encodes both equivalence and timing preferences into cut weights, adaptively updates the weights and re-searches the cut according to six failure modes, and embeds estimated wire load into the inner-loop acceptance criterion to filter physically robust candidates via two-stage gating (an ideal-net improvement threshold and a simplified SPEF recheck). Strict timing improvements are achieved on ISCAS89 (8/8), ITC-99 (18/19), and PicoRV32 (2/3) benchmarks. Under the convergence configuration, the average WNS improvement is 1.07 ns (mean of per-circuit improvements), 6.7$\times$ that of the gate-sizing-only method; under the efficiency-oriented configuration, only 27 STA calls are needed to complete the 8/8 repair, cutting the overhead by 99.1\%. After real placement and routing, 5 of the 8 circuits retain the improvement, 1 ties, and 2 regress by no more than 0.02 ns. FAECO implements a closed-loop failure-driven logic re-synthesis repair on the open-source toolchain.}\par
\vspace{3pt}
{\bfseries Key words:} {\small timing ECO; re-synthesis; bi-objective cut; failure-aware feedback; simplified SPEF gating; pre-placement logic re-synthesis}\par
\end{flushleft}
]
Body text.
\end{document}